\Article{12}
Es bezeichnen \(x,x',x''\) etc. unbestimmt die voneinander unabhängigen Fehler von Beobachtungen derselben Art, deren relative Wahrscheinlichkeiten durch die vorangestellte Charakteristik \(\varphi\) ausgedrückt werden; ferner sei \(y\) eine gegebene rationale Funktion der Unbestimmten \(x,x',x''\) etc.: Das mehrfache Integral (I)
\[
\int\varphi x.\varphi x'.\varphi x''\mathord{.}\,\mathord{.}\,\mathord{.}\,\mathord{.}\,\mathord{.}\,\mathrm{d}x.\mathrm{d}x'.\mathrm{d}x''\mathord{.}\,\mathord{.}\,\mathord{.}\,\mathord{.}
\]
über alle Werte der Unbestimmten \(x,x',x''\) erstreckt, für die der Wert von \(y\) innerhalb der gegebenen Grenzen \(0\) und \(\eta\) liegt, drückt die Wahrscheinlichkeit eines Wertes von \(y\) aus, der unbestimmt innerhalb \(0\) und \(\eta\) liegt. Offensichtlich wird dieses Integral eine Funktion von \(\eta\) sein, deren Differential wir \(=\psi\eta.\mathrm{d}\eta\) setzen wollen, sodass das Integral selbst dem Integral \(\int\psi\eta.\mathrm{d}\eta\) von \(\eta=0\) an gleich wird. Auf diese Weise ist zugleich die Charakteristik \(\psi\eta\) als Ausdruck der relativen Wahrscheinlichkeit jedes Wertes von \(y\) anzusehen. Da \(x\) als Funktion der Unbestimmten \(y,x',x''\) etc. betrachtet werden kann, die wir
\[
=f(y,x',x''\mathord{.}\,\mathord{.}\,\mathord{.}\,\mathord{.})
\]
setzen wollen, wird das Integral (I)
\[
\begin{aligned}
=\int\varphi.f(y,x',x''\mathord{.}\,\mathord{.}\,\mathord{.}).\frac{\mathrm{d}f(y,x',x''\mathord{.}\,\mathord{.}\,\mathord{.}\,\mathord{.})}{\mathrm{d}y}.{}\\
\varphi x'.\varphi x''\mathord{.}\,\mathord{.}\,\mathord{.}\,\mathord{.}\,\mathrm{d}y.\mathrm{d}x'.\mathrm{d}x''\mathord{.}\,\mathord{.}\,\mathord{.}\,\mathord{.}
\end{aligned}
\]
wobei \(y\) von \(y=0\) bis \(y=\eta\) zu erstrecken ist, die übrigen Unbestimmten aber über alle Werte, denen ein reeller Wert von \(f(y,x',x''\mathord{.}\,\mathord{.}\,\mathord{.}\,\mathord{.})\) entspricht. Hieraus \SourceLeaf{018}folgt
\[
\begin{aligned}
\psi y=\int\varphi.f(y,x',x''\mathord{.}\,\mathord{.}\,\mathord{.}\,\mathord{.}).\frac{\mathrm{d}f(y,x',x''\mathord{.}\,\mathord{.}\,\mathord{.}\,\mathord{.})}{\mathrm{d}y}.{}\\
\varphi x'.\varphi x''\mathord{.}\,\mathord{.}\,\mathord{.}\,\mathord{.}\,\mathrm{d}x'.\mathrm{d}x''\mathord{.}\,\mathord{.}\,\mathord{.}\,\mathord{.}
\end{aligned}
\]
wobei die Integration, bei der \(y\) als konstant betrachtet werden muss, über alle Werte der Unbestimmten \(x',x''\) etc. erstreckt wird, die \(f(y,x',x''\mathord{.}\,\mathord{.}\,\mathord{.}\,\mathord{.})\) einen reellen Wert geben.

\Article{13}
Um diese Integration tatsächlich auszuführen, wäre die Kenntnis der Funktion \(\varphi\) erforderlich, die meistens unbekannt ist: Ja selbst wenn diese Funktion bekannt wäre, würde die Integration in den meisten Fällen die Kräfte der Analysis übersteigen. Unter diesen Umständen werden wir zwar die Wahrscheinlichkeit der einzelnen Werte von \(y\) nicht angeben können: Anders verhält es sich aber, wenn nur der Mittelwert von \(y\) gesucht wird, der aus der Integration von \(\int y\psi y.\mathrm{d}y\) über alle Werte von \(y\) entsteht, die es erreichen kann. Da nun offensichtlich für alle Werte, die \(y\) nicht erreichen kann, sei es wegen der Natur der Funktion, die es ausdrückt (z. B. für negative Werte, wenn \(y=xx+x'x'+x''x''+\text{etc.}\) wäre), sei es, weil den Fehlern \(x,x',x''\) etc. selbst bestimmte Grenzen gesetzt sind, \(\psi y=0\) zu setzen ist, bleibt das Ergebnis offensichtlich dasselbe, wenn jene Integration über alle reellen Werte von \(y\) erstreckt wird, nämlich von \(y=-\infty\) bis \(y=+\infty\). Nun ist das Integral \(\int y\psi y.\mathrm{d}y\) zwischen bestimmten Grenzen, nämlich von \(y=\eta\) bis \(y=\eta'\) genommen, gleich dem Integral
\[
\begin{aligned}
\int y\varphi.f(y,x',x''\mathord{.}\,\mathord{.}\,\mathord{.}\,\mathord{.})\frac{\mathrm{d}f(y,x',x''\mathord{.}\,\mathord{.}\,\mathord{.}\,\mathord{.})}{\mathrm{d}y}.{}\\
\varphi x'.\varphi x''\mathord{.}\,\mathord{.}\,\mathord{.}\,\mathord{.}\,\mathrm{d}y.\mathrm{d}x'.\mathrm{d}x''\mathord{.}\,\mathord{.}\,\mathord{.}\,\mathord{.}
\end{aligned}
\]
wobei die Integration von \(y=\eta\) bis \(y=\eta'\) sowie über alle Werte der Unbestimmten \(x',x''\) etc. erstreckt wird, denen ein reeller Wert von \(f(y,x',x''\mathord{.}\,\mathord{.}\,\mathord{.}\,\mathord{.})\) entspricht; oder, was dasselbe ist, gleich dem Wert des Integrals
\[
\int y\varphi x.\varphi x'.\varphi x''\mathord{.}\,\mathord{.}\,\mathord{.}\,\mathord{.}\,\mathrm{d}x.\mathrm{d}x'.\mathrm{d}x''\mathord{.}\,\mathord{.}\,\mathord{.}\,\mathord{.}
\]
wobei in dieser Integration für \(y\) sein durch \(x,x',x''\) etc. ausgedrückter Wert einzusetzen und die Integration über alle Werte dieser Unbestimmten zu erstrecken ist, denen ein Wert von \(y\) zwischen \(\eta\) und \(\eta'\) entspricht. Hieraus schließen wir, dass das Integral \(\int y\psi y.\mathrm{d}y\) über alle Werte von \(y\), von \(y=-\infty\) bis \(y=+\infty\) erstreckt, durch die Integration von
\[
\int y\varphi x.\varphi x'.\varphi x''\mathord{.}\,\mathord{.}\,\mathord{.}\,\mathord{.}\,\mathrm{d}x.\mathrm{d}x'.\mathrm{d}x''\mathord{.}\,\mathord{.}\,\mathord{.}\,\mathord{.}
\]
\SourceLeaf{019}über alle reellen Werte von \(x,x',x''\) etc. erhalten wird, nämlich von \(x=-\infty\) bis \(x=+\infty\), von \(x'=-\infty\) bis \(x'=+\infty\) etc.

\clearpage
